% Transcription — fonds Alexandre Grothendieck, shelfmark 18, batch 3 (pages 41-60).
% Established from the facsimile archives/batches/18/batch-03.pdf (a mirror of
% https://grothendieck.umontpellier.fr/18.pdf), per the skill
% transcribe-grothendieck. Its PDF page 1 is archivists' page 41 (checked
% against the pencilled number bottom-left, and again at 50, 52, 54, 56, 58, 60);
% the batch file has no cover sheet.
%
% Pass: Opus 5 (claude-opus-5), 2026-09-16 — first pass, unchecked against the pages by a human.
%
% Eleven leaves carry his hand: 41, 42, 44, 46, 48, 50, 52, 54, 56, 58, 60 — the
% even ones and the cover. The odd ones, 43 45 47 49 51 53 55 57 59, are the
% backs of typescripts foreign to these notes (an English mathematical
% typescript, Bourbaki fascicles) and are skipped: the gaps in the numbering are
% the record. Page 41 is otherwise blank and carries only the two lines that
% name the run; per the skill it takes no \page and its words open the file as
% the section heading.
%
% Batches 1, 2 and 4 are being transcribed at the same time by other passes, so
% no earlier .tex fixed the notation. No run crosses the 40/41 boundary: page 40
% is on Hodge theory and a real subfield for T, page 41 opens a fresh heading and
% page 42 a fresh list of problems. A run DOES cross the far boundary: page 60
% breaks off on « ... est un sous-groupe ouvert de G(A_C) » and page 61 takes the
% same computation up, with the relation M^Z_w = g M^Z_v and the conjugated
% cocycle phi_w(sigma) = g_infty phi_v(sigma) g_infty^{-1}. Page 61 is batch 4's;
% only 41-60 are transcribed here, and a \note marks the break.
%
% What the batch is about. A transcendence programme for periods, in the
% language of motives and their Mumford-Tate groups. Page 42 lists four problems
% of comparison (Betti/de Rham lattices, Betti lattices for different
% topologies, real structures and splittings, de Rham against p-adic over a
% complete valued field, and a p-adic cohomology at every p). Page 44 works the
% example C x C' of two elliptic curves without complex multiplication: G of
% dimension 7, B of dimension 5, and a transcendence degree bounded by 5 where
% dim G - 1 = 6 was expected. Pages 46-50 are the general argument: for a motive
% M over k, each place v gives a lattice M^Z_v inside the product of the M(l),
% the places v_1, ..., v_n give principal homogeneous spaces P_i = Isom(xi,
% xi_{v_i}) under the Mumford-Tate group G, and a point of the double coset
% space bounds the transcendence degree by (n-1)d, d = dim G. A lemma on generic
% elements of a coset U.a in G(Q-hat), and the strengthened Tate conjecture
% (Gal(k/k_0) -> G(Z-hat) of finite index), let the v_i be moved by Galois until
% the components are independently generic. Pages 52-58 turn to the archimedean
% place: with k embedded in C, the adelic module M^{A_C} (l finite OR infinite),
% the semi-linear map sigma_infty coming from M_k(infty) = M_{k_0}(infty) tensor
% k, a map phi_v : Gal(k/k_0) -> G(A_C)/G(Z-hat), then — the point he marks with
% three exclamation marks — the archimedean component alone is well determined
% in G(k), giving phi_{v,infty} : Gal(k/k_0) -> G(k) inside GL(M_k(infty)), a
% 1-cocycle, hence a principal homogeneous space P_{v,k} under G with Galois
% descent data and a form P_{v,k_0} over k_0. Page 60 asks that the image be
% Zariski-dense, resp. open in G(A_C).
%
% Notation fixed for the batch: M for a motive, M(l) for its l-adic realisation,
% M^Z_v for the lattice attached to a place v, M^{A_C} for the adelic module;
% A_C for his adeles with C at the infinite place (the superscript on page 52
% and the argument of G on pages 54, 58 and 60 are the same letter, and page 60
% writes it as an unambiguous capital C); Q-hat and Z-hat as \widehat{\mathbb{Q}}
% and \widehat{\mathbb{Z}}; G = G_{M,xi_0} the Mumford-Tate group, G(Z_l),
% G(Z-hat), G(k); sigma-dot and sigma_* for the operator of a Galois element,
% sigma_infty for its archimedean (semi-linear) component; phi_v and psi_v for
% the cocycles; k_0 the field of definition, k its extension embedded in C.
% « 1/2 linéaire » is transcribed « semi-linéaire », his 1/2 for semi-. His « ré »
% is drawn so that « réseaux » reads « viscaux » and « réelles » reads
% « vielles »; the readings are fixed once here and held.
\documentclass[11pt,a4paper]{article}
% Le préambule commun à toutes les transcriptions du fonds Grothendieck.
%
% Il tient en une idée : **l'incertitude se compose, elle ne se résout pas.**
% Une transcription qui lisse un mot douteux a détruit la seule chose pour
% laquelle on la faisait — un lecteur doit pouvoir savoir, à chaque endroit,
% ce qui était sur le feuillet et ce qui a été ajouté. Les six macros qui
% suivent sont donc l'appareil critique, et non de la décoration.
%
% Les mêmes commandes sont reconnues par `scripts/render.mjs`, qui produit la
% vue de lecture du site. Une macro ajoutée ici doit l'être aussi là-bas,
% faute de quoi le rendu échouera bruyamment — ce qui est voulu : un rendu
% silencieusement faux est pire qu'un rendu absent.

\ProvidesPackage{grothendieck}[2026/08/08 Transcriptions du fonds Grothendieck]

\usepackage{amsmath,amssymb,amsthm}
\usepackage{mathrsfs}
\usepackage{stmaryrd}
% Les diagrammes commutatifs sont partout dans ce fonds : c'est la langue
% même de la géométrie algébrique de Grothendieck, et une transcription qui
% les rendrait en prose ne transcrirait rien.
\usepackage{tikz-cd}
% Français par défaut — c'est la langue des feuillets — mais l'anglais est
% chargé aussi : la traduction et le résumé sont des documents anglais, et la
% typographie française (espaces avant les deux-points) y serait une faute.
% Ces documents basculent par \selectlanguage{english} en tête de corps.
\usepackage[english,french]{babel}
\usepackage{fontspec}
\usepackage{geometry}
\geometry{a4paper,margin=25mm,marginparwidth=28mm}
\usepackage{marginnote}
\usepackage{xcolor}
% ulem, not soul. `soul`'s \st and \ul are built on a character-by-character
% scan that XeLaTeX's font machinery breaks: the document fails with an
% undefined control sequence rather than degrading. ulem does the same two jobs
% and is Unicode-engine safe. `normalem` stops it hijacking \emph, which the
% transcriptions use for Grothendieck's own emphasis.
\usepackage[normalem]{ulem}
\usepackage{hyperref}
\hypersetup{colorlinks=true,linkcolor=black,urlcolor=black,pdfborder={0 0 0}}

% --- Métadonnées du lot -----------------------------------------------------
% Reprises telles quelles par le script de rendu, qui les affiche en tête de
% la vue de lecture. Elles fixent ce que le fichier prétend couvrir : sans
% elles, une transcription flotte sans point d'attache dans un fonds de seize
% mille feuillets.

\newcommand{\folder}[1]{\gdef\grothFolder{#1}}
\newcommand{\batch}[1]{\gdef\grothBatch{#1}}
\newcommand{\leaves}[2]{\gdef\grothFirst{#1}\gdef\grothLast{#2}}
\newcommand{\foldertitle}[1]{\gdef\grothFoldertitle{#1}}
\gdef\grothFolder{?}\gdef\grothBatch{?}\gdef\grothFirst{?}\gdef\grothLast{?}\gdef\grothFoldertitle{}

% --- Le résumé --------------------------------------------------------------
%
% Il ouvre la lecture modernisée, sous le titre « Résumé », et il est composé
% en petit. Ce n'est pas un ornement : c'est le seul passage du document qui
% ne soit pas des mathématiques, et un lecteur doit pouvoir le reconnaître
% avant de l'avoir lu — puis le sauter s'il connaît déjà le sujet, ou s'y
% arrêter s'il ne le connaît pas. Le filet qui le suit marque la reprise du
% corps.

\newenvironment{resume}
  {\small\setlength{\parskip}{0.45em}}
  {\par\medskip\hrule\medskip}

% --- Mots-clés ---------------------------------------------------------------
%
% En anglais, en clôture du résumé de la lecture modernisée : le vocabulaire
% moderne sous lequel chercher ce que les feuillets construisent. Le manifeste
% du site (`npm run manifest`) les extrait pour étiqueter la cote entière —
% cette ligne est la source unique des tags, il n'y a pas de fichier à côté.
\newcommand{\keywords}[1]{\par\smallskip\noindent
  {\footnotesize\textsc{Keywords} --- \textit{#1}}\par}

% --- L'appareil critique ----------------------------------------------------

% Le numéro de feuillet, en marge. C'est l'ancre qui permet de régler un
% désaccord sur une formule en regardant *un* feuillet plutôt qu'une chemise
% de six cents. Le site s'en sert aussi pour tourner les pages du fac-similé
% quand on fait défiler la transcription.
\newcommand{\leaf}[1]{%
  \marginnote{\footnotesize\color{gray!70}\textsf{#1}}%
  \label{leaf:#1}\ignorespaces}

% Illisible : on ne devine pas. Un mot inventé qui se lit comme les autres est
% le pire résultat possible.
\newcommand{\ill}{\textcolor{red!60!black}{[\,\ldots\,]}}

% Lecture douteuse : le mot est donné, et signalé comme incertain.
\newcommand{\uncertain}[1]{\uline{#1}}

% Ajout de l'éditeur — une lettre manquante, un mot sous-entendu.
\newcommand{\add}[1]{\textcolor{blue!50!black}{[#1]}}

% Biffé par Grothendieck lui-même. Ce qu'il a rayé fait partie du document :
% une rature dit souvent où la pensée a changé de direction.
\newcommand{\struck}[1]{\sout{#1}}

% Note du transcripteur, sur l'état matériel du feuillet.
\newcommand{\note}[1]{\par\smallskip{\small\color{gray!60!black}\textit{[#1]}}\par\smallskip}

% Note marginale de Grothendieck, distincte du corps du texte.
\newcommand{\marginal}[1]{\par\smallskip\noindent\hspace{1em}%
  {\small\color{gray!60!black}#1}\par\smallskip}

% --- Titre ------------------------------------------------------------------

\AtBeginDocument{%
  \begin{center}
    {\large\bfseries Fonds Alexandre Grothendieck --- cote n\textsuperscript{o}~\grothFolder}\\[2pt]
    {\itshape\grothFoldertitle}\\[4pt]
    {\small feuillets \grothFirst--\grothLast\ (lot \grothBatch)}\\[2pt]
    {\footnotesize\color{gray!60!black}Transcription établie d'après le fac-similé numérisé
     par l'Université de Montpellier.\\
     Les lectures incertaines sont soulignées, les illisibles notées [\,\ldots\,],
     les ajouts entre crochets.}
  \end{center}
  \vspace{1em}}

\endinput


\folder{18}
\batch{3}
\pages{41}{60}
\dating{[à partir de 1969-vers 1972]}
\watermark{Édition de démonstration}
\foldertitle{Motifs et théorie de Hodge : notes manuscrites (s.d.)}

\begin{document}

\section*{Variations du réseau discret en termes d'une place archimédienne. Critère de transcendance}

\note{Le titre est de lui. Ces deux lignes occupent seules le haut du feuillet
41, par ailleurs vierge : ce feuillet sert de couverture à ce qui suit et n'est
pas transcrit comme page, d'où le saut de 41 à 42 dans la numérotation.
« place » y est abrégé « pl. ». Les autres feuillets absents (43, 45, 47, 49,
51, 53, 55, 57, 59) portent au dos des textes étrangers à ces notes.}

\page{42}

\emph{Pbs}

\begin{enumerate}
  \item[1)] Betti, DeRham
  \item[a)] Comparaison « réseaux » de DeRham et \struck{de} de Betti
  \item[b)] Comparaison réseaux de Betti pour topologies différentes ?
  \item[c)] Comparaison structures réelles, resp.\ splittages, (dans DeRham)
        pour top.\ diff.
  \item[2)] Betti, $\ell$-adique. \struck{Comp} Comparaison de réseaux de Betti
        pour top.\ diff.\ (dans ch.\ $\ell$-adique) --- comp.\ \uncertain{évident ou $p$}
  \item[3)] DeRham, $p$-adique. Sur un corps valué (\uncertain{complet} \ill{})
        (non archimédien), comparaison possible entre cohomologie $p$-adique et
        cohomologie de DeRham ?
  \item[4)] Définition de cohomologie $p$-adique en tout $p$ ??
\end{enumerate}

\note{Le mot entre guillemets au point a) est le même que celui du feuillet de
couverture : sa « ré » se lit « vi », et « réelles » du point c) se lit de même
« vielles ». Deux insertions interlinéaires : au-dessus de « corps valué » du
point 3), trois mots dont seul le premier se laisse lire ; au-dessus du point
2), « (dans ch.\ $\ell$-adique) », donné ici dans la ligne.}

\page{44}

$C$, $C'$ courbes elliptiques sans mult.\ complexe, avec \ill{}
\note{Les derniers mots de la ligne sont comprimés contre le bord droit et ne se
laissent pas lire ; ce qui suit ($G \simeq SL(2) \times SL(2) \times
\mathbb{G}_{m}$, de dimension $7$) demande que $C$ et $C'$ soient non
isogènes.}

$C \times C' = A$

\struck{isogénie} $G \simeq SL(2) \times SL(2) \times \mathbb{G}_{m}$ \quad de dim.\ $7$

$B \simeq B(2) \times B(2) \times \mathbb{G}_{m}$ \quad de dim.\ $5$

Alors $K$ de degré de transcendance $\leqslant 5$ sur $k$, donc de degré de
transcendance $\leqslant 5$ sur $k(\pi)$ [alors qu'on s'attend à
$\dim G - 1 = 6$].

\marginal{NB Ceci \uncertain{joint} pour \emph{ttes} incarnations $k \subset \mathbb{C}$ !}

\page{46}

$M$ motif sur $k$ \add{alg.\ clos}, $k$ isomorphe à un sous-corps de
$\mathbb{C}$.

Pour toute place $v$ de $k$, on a $M^{\mathbb{Z}}_{v}$ qui est un sous-réseau de
$\prod_{\ell\ \mathrm{fini}} M(\ell)$. On veut établir les relations entre les
$M^{\mathbb{Z}}_{v}$ pour diverses valeurs $v_{1}, \ldots, v_{n}$ de $v$.

Choisissons un \uncertain{sous-foncteur} de \uncertain{référence} (p.\ ex.\
défini sur $\mathbb{Q}$)
\[
  \xi : M \rightsquigarrow M_{\xi_{0}} \otimes_{\xi_{0}} \quad (= T(\mathbb{Q}))
\]
avec des isom.\ foncto\supplied{riels} $\xi_{\struck{\ill{}}}(M)_{\mathbb{Q}_{\ell}}
= M(\ell) \otimes_{\mathbb{Z}_{\ell}} \mathbb{Q}_{\ell}$.
\note{L'indice de $\otimes$ dans la première formule et la source du second
membre de la seconde sont l'un et l'autre douteux ; on donne la lecture que
demande la suite.}

Alors $v_{1}, \ldots, v_{n}$ définissent les fibrés principaux homogènes
\[
  P_{i} = \underline{\mathrm{Isom}}(\xi, \xi_{v_{i}})
\]
sous $G^{\xi}_{M,\xi_{0}}$ (ou sur $\mathbb{Q}$), et \ill{} un point dans
$\bigl(\prod P_{i}\bigr)(\widehat{\mathbb{Q}})$, où $\widehat{\mathbb{Q}} =
\prod_{\ell\ \mathrm{fini}} \mathbb{Q}_{\ell}$, d'où un pt de
$\bigl[\bigl(\prod P_{i}\bigr)/G\bigr](\widehat{\mathbb{Q}})$.

On peut aussi considérer que l'on a un pt canonique dans
\struck{$GL(n) \times \cdots \times GL(n)$}
\struck{$\bigl[GL(n) \times \cdots \times GL(n)\bigr](\widehat{\mathbb{Q}}) / GL(n)(\widehat{\mathbb{Q}})$}
\[
  GL(n, \widehat{\mathbb{Z}}) \times \cdots \times GL(n, \widehat{\mathbb{Z}})
  \,\backslash\, GL(n, \widehat{\mathbb{Q}}) \times \cdots \times GL(n, \widehat{\mathbb{Q}})
  \,/\, GL(n, \widehat{\mathbb{Q}})
\]
\add{$n = \mathrm{rang}\ M$}
\note{Dans la ligne encadrée il avait d'abord écrit $\mathbb{Z}$, biffé et
récrit $\widehat{\mathbb{Z}}$ ; le $GL$ du second membre biffé est lui-même
surchargé. En marge gauche, en face de ce bloc, deux mots qui ne se lisent pas.}

Les considérations qui précèdent montrent que le degré de transcendance des pts
en question

\page{48}

est majoré par $(n-1)d$, car \ill{} $\bigl((n-1) \cdot n\bigr)$ où
$d = \dim G$. \struck{\ill{}}

D'ailleurs, si $v_{1}, v_{2}, \ldots, v_{n}$ sont donnés, \struck{\ill{}}
définissant \ill{} $g_{1}, \ldots, g_{n} \in G(\widehat{\mathbb{Z}})$, en
prenant $\tau_{2}, \ldots, \tau_{n} \in \mathrm{Gal}(k/k_{0})$ (si $k_{0}$ est
\ill{} corps de \struck{définition} \add{de type fini} de $M$) et $v'_{1},
v'_{2}, \ldots, v'_{n}$ avec $v'_{i} = \tau_{i} v_{i}$, $1 \leqslant i
\leqslant n$, on trouvera les éléments \struck{$g$} $\dot\tau_{2} g_{2}, \ldots,
\dot\tau_{n} g_{n}$, \ill{} $\dot\tau_{i}$ l'image de \ill{} dans
$G(\widehat{\mathbb{Z}})$. \struck{\ill{}}

\marginal{N.B. On \ill{} \uncertain{prendra} \ill{} la « fine » \ill{} donc conne\ill{}}

Admettant la \emph{conjecture de Tate renforcée}, $\mathrm{Gal}(k/k_{0}) \to
G(\widehat{\mathbb{Q}})$ a une image \uncertain{$\sigma$} qui est d'indice fini,
il s'ensuit que quitte à remplacer les $v_{i}$ par des $\tau_{i} v_{i}$
convenables, on peut \struck{obtenir} $g_{2}, \ldots, g_{n}$ génériques
indépendamment pour \emph{tt} $\ell$, si \ill{} devrait ceci :

\begin{quote}
Soit $G$ un gr.\ alg.\ \struck{\ill{}} affine connexe sur $\mathbb{Q}$, $U$
s\supplied{ous}-g\supplied{rou}pe \ill{} ouvert de $G(\widehat{\mathbb{Q}})$, $P = U.a$
\add{une classe} dans $G(\widehat{\mathbb{Q}})$ mod $U$. Alors il existe dans
$P$ un élt \ill{} dont toutes les composantes soient génériques rel.\
$\mathbb{Q}$.
\end{quote}

\note{Le lemme est enfermé dans un grand crochet, en bas à droite du feuillet ;
il est repris à la première ligne de la page 50.}

\page{50}

Donc en admettant ce lemme, on trouvera que \emph{tt} syst.\ $v_{1}, \ldots,
v_{n}$ \ill{} générique, \ill{} pt \struck{générique} de $V(\widehat{\mathbb{Z}})$
\struck{générique} tel que pour \emph{tt} $\ell$, \ill{} $V(\mathbb{Z}_{\ell})$
soit générique rel.\ à $\mathbb{Q}$.
\note{Le $V$ de ces trois lignes n'est pas le $G$ du lemme ; la lettre est
écrite en deux traits très ouverts et le reste de la page n'y revient pas.}

D'ailleurs, on \ill{} inversement, \ill{}

\marginal{Moyennant \ill{} complétion \ill{} le \ill{} de $M$, \ill{} $G$ \ill{}
$T$ \ill{} qui \ill{} \ill{} des \ill{} $M^{\mathbb{Z}}_{v}$ \ill{}
$(M_{f} = \prod \mathbb{Z}_{\ell})$}
\note{Cette note est écrite de biais dans la marge gauche, en travers des
lignes ; un lambeau en est entouré à la mine (« le lemme \ill{} »). On ne donne
que le peu qui se laisse lire.}

$M^{\mathbb{Z}}_{v}$ \ill{} $v_{i}$ \ill{} $M^{\mathbb{Z}}_{v_{i}}$ \ill{}
comme \ill{}

\ill{} \ill{} en transformé des réseaux \ill{} comparables.

[On \ill{} que le choix des $\tau_{i}$, et enfin $v_{i}$,
\emph{dépend de $M$} ; \ill{} que pour $M$ variable, \ill{} bien déterminé les
\ill{} par des \ill{} $\xi$.]

\page{52}

Supposons maintenant $k$ isomorphe à $\mathbb{C}$, et donnons-nous une \ill{}
pour laquelle \ill{} \emph{complet}. Alors on a \ill{} les réseaux
\[
  M^{\mathbb{Z}}_{v} \subset M^{\mathbb{A}_{\mathbb{C}}}
  \quad\text{où}\quad
  M^{\mathbb{A}_{\mathbb{C}}} = \struck{\ill{}} \prod M(\ell)
\]
\marginal{$\ell$ fini \emph{ou infini}}
\ill{} indépendant.

\struck{Comparaison \ill{}} Si $M$ est défini sur \ill{} $k_{0}$, et $k_{0}$
\ill{} pour \ill{} que $G_{M}$ soit connexe \ill{} \ill{} de transcendance
\struck{\ill{}} \ill{} minimum, \ill{} alors \struck{\ill{}} on a, pour
$\sigma \in \struck{\mathrm{Aut}_{k_{0}}\ \ill{}}\ \mathrm{Gal}(k/k_{0})$,
\[
  M^{\mathbb{Z}}_{\sigma v} = \dot\sigma M^{\mathbb{Z}}_{v} ,
\]
\marginal{NB $k$ n'est pas \struck{\ill{}} \ill{} algébrique sur $k_{0}$, mais
de degré de transcendance \ill{}}
où $\dot\sigma$ opère sur $M^{\mathbb{A}_{\mathbb{C}}}$ \ill{} \emph{les
images} \ill{} :
\[
  \dot\sigma\bigl((x_{\ell})_{\ell\ \mathrm{fini}},\ x_{\infty}\bigr)
  = \bigl(\tau_{\ell} x_{\ell},\ \dot\sigma_{\infty} x_{\infty}\bigr)
\]
où $\tau_{\ell}$ provient de $\mathrm{Gal}(k/k_{0}) \to G(\mathbb{Z}_{\ell})$,
et $\dot\sigma_{\infty}$ est l'hom.\ \emph{semi-}linéaire déduit de $\sigma$ via
$M_{k}(\infty) \simeq M_{k_{0}}(\infty) \otimes_{k_{0}} k$.
\note{Il écrit « $\tfrac{1}{2}$ linéaire » : son $\tfrac{1}{2}$ pour
\emph{semi-}, comme ailleurs « $\tfrac{1}{2}$cosimplicial ». Le $\mathbb{A}$ de
$M^{\mathbb{A}_{\mathbb{C}}}$ porte ici un petit c ; la page 60 écrit le même
groupe $G(\mathbb{A}_{\mathbb{C}})$ avec un C capital non ambigu, et le facteur
archimédien est bien $G(\mathbb{C})$.}

\page{54}

Comme \ill{} $v$ et $w$, \ill{} transformé dans \struck{Dont} comme d'habitude
\ill{} il existe un
\[
  g \in G(\mathbb{A}_{\mathbb{C}}) = \prod_{\ell\ \mathrm{fini}} G(\mathbb{Z}_{\ell}) \times G(\mathbb{C})
\]
\struck{dit que $M^{\mathbb{Z}}_{w} = g$} tel que
\[
  \underline{M}^{\mathbb{Z}}_{w} = g\, M^{\mathbb{Z}}_{v}
\]
\add{où $G = G_{M,\xi_{0}}$} \ill{} isomorphes, \ill{} pour \ill{}, mod
multiplication $g \rightsquigarrow gh$ ($h \in G(\widehat{\mathbb{Z}})$).

$g$ est \ill{} associé à $M$, $v$. \ill{} $N_{v}$ et $N_{w}$ \ill{}

D'ailleurs, $g$ est déterminé \ill{} $h \in G(\widehat{\mathbb{Z}})$.

\marginal{N.B. \ill{} de \ill{} fixe $g_{\sigma}$ \ill{} indéterminé mod
$G(\widehat{\mathbb{Z}})$ \ill{} $g_{\ell}$ \ill{} déterminé}

\struck{$M^{\mathbb{Z}}_{\sigma v} = \dot\sigma_{\infty}(g M^{\mathbb{Z}}_{v})$}
\[
  M^{\mathbb{Z}}_{\sigma v} = \struck{g_{\sigma}}\ g_{\sigma} M^{\mathbb{Z}}_{v}
  = \dot\sigma_{\infty} M^{\mathbb{Z}}_{v}
\]
d'où \ill{} \emph{une} application
\[
  \mathrm{Gal}(k/k_{0}) \xrightarrow{\ \varphi_{v}\ } G(\mathbb{A}_{\mathbb{C}})/G(\widehat{\mathbb{Z}})
\]
\[
  M^{\mathbb{Z}}_{\sigma v} = g_{\sigma} \varphi_{v}(\sigma) M^{\mathbb{Z}}_{v}
  \quad\text{i.e.}\quad
  \dot\sigma_{\infty}(M^{\mathbb{Z}}_{v}) = \varphi_{v}(\sigma) M^{\mathbb{Z}}_{v}
\]
\marginal{N.B. $\sigma \rightsquigarrow$ \ill{} une surjection \ill{} peut
induire \ill{} $v|k_{0}$ \ill{} $k_{0}$, \ill{} de card.\ $2$ (conj.\ complexe)}

\page{56}

\struck{Fixons une $v$ de \ill{} de $k$, \ill{}}

On \ill{} $g_{\sigma}$ est \ill{} co\ill{} par $G(\widehat{\mathbb{Z}})$,
d'\ill{} (où \ill{} $\ell$ = indéterminé) \ill{} l'application \ill{} par
prendre
\[
  g_{\sigma} = \bigl((\sigma_{\ell})_{\ell\ \mathrm{fini}},\ \dot\sigma_{\infty}\bigr)
\]
\ill{} maintenant $\dot\sigma_{\infty}$ est \ill{} \emph{bien déterminé} dans
$G(k)$ !!!

On obtient ainsi une application canonique
\[
  \mathrm{Gal}(k/k_{0}) \xrightarrow{\ \varphi_{v,\infty}\ } G_{\struck{\ill{}}}(k)
  \subset GL\bigl(M_{k}(\infty)\bigr)
\]
\marginal{N.B. $G$ peut ici être considéré comme \ill{} défini sur $k_{0}$,
\ill{} indépendant de $v$ !}

\struck{et telle qui \ill{} ait} $\dot\sigma M^{\mathbb{Z}}_{v} = \sigma_{*}
M^{\mathbb{Z}}_{v}$ pour $\sigma \in \mathrm{Gal}(k/k_{0})$, avec
\[
  \sigma_{*} = \bigl((\sigma_{\ell})_{\ell\ \mathrm{fini}},\ \varphi_{v}(\sigma)\bigr) .
\]

On vérifie aussitôt, i.e.\ posant $\psi_{v} = \varphi_{v}^{-1}$ :
\[
  \varphi_{v}(\tau\sigma) = {}^{\tau}\varphi_{v}(\sigma)\, \varphi_{v}(\tau) ,
  \qquad
  \psi_{v}(\tau\sigma) = \psi_{v}(\tau)\, {}^{\tau}\psi_{v}(\sigma) .
\]
\note{Les deux formules portent chacune un astérisque biffé en bout de ligne ;
le premier membre de la seconde a d'abord été écrit $\varphi$, corrigé en
$\psi$.}

\page{58}

Ainsi $\varphi$ est un \emph{$1$-cocycle}, quand on \struck{fait} opère \ill{}
$\mathrm{Gal}(k/k_{0})$ sur $G(k)$ de $2$ façons évidentes \ill{} (par
conjugaison \ill{}). Il correspond à un fibré principal \ill{} homogène
$P_{v,k}$ sous $G(k)$, avec opérations de $\mathrm{Gal}(k/k_{0})$ sur
$\bigl(P_{v}, G(k)\bigr)$ \struck{compatible \ill{}} \ill{} fibré principal
\ill{} $G_{k_{0}}$, avec \ill{}

\ill{} points de $\mathrm{Gal}(k/k_{0})$ \ill{} \struck{compatible}

Je désigne $P_{v,k_{0}}$ \ill{} le \ill{} \ill{} \uncertain{splitt.}\ \ill{}
donné sur $k$.

D'ailleurs, \ill{} c'est le \ill{} fibré principal $\xi_{M,\infty}$ avec
$\xi^{\mathbb{Z}}_{M,v} \otimes_{\mathbb{Z}} k_{0}$, \ill{} en effet obtenu
\ill{} de transcendance (N\textsuperscript{o} 1, \struck{cette fini} \ill{})
\ill{} $P_{v}^{k}$ \ill{} l'\struck{multiplicité} multiplicité \ill{} dire que
\[
  \mathrm{Gal}(k/k_{0}) \xrightarrow{\ \varphi_{v}\ } G(k)
  \quad\bigl[\subset GL(M_{k}(\infty))\bigr]
\]
\ill{} rang de Zariski \ill{} dense.

\marginal{À vérifier / \emph{ok}.}
\marginal{$\pi$ transporte \ill{} au $k_{0}$\ill{}}
\note{Les deux notes précédentes sont écrites dans la marge gauche, la première
barrée d'un grand trait oblique.}

\page{60}

\ill{} (pour \ill{} : \ill{}) $k_{0}$, $\sigma \rightsquigarrow
({}^{\sigma}g) g^{-1}$ \add{\ill{} image de $g^{-1}g^{\sigma}$} une extension
finie de \struck{[n'est déterminé \ill{} pour $g \in G(k_{0})$]} générique
$/k_{0}$. Car \ill{} $g \in G(k)$ \ill{} générique sur $G(k)$, et \ill{} pour
$g^{\sigma}$ \emph{tt} \ill{}

\ill{} On voit \ill{} que $g^{\sigma} g$ \ill{} peut \struck{tt} être
\uncertain{exact} \ill{} On voit \ill{} que $g$ soit \uncertain{générique}
$/k_{0}$, \ill{} telle que \ill{} \struck{($\sigma$)} \ill{} $k_{0}(g)$.
\struck{Et d'autres}

\ill{} $\varphi$ a \ill{} génér\ill{} \struck{\ill{}}

Considérons $\sigma \rightsquigarrow \sigma_{*}$,
\[
  \mathrm{Gal}(k/k_{0}) \longrightarrow G(\mathbb{A}_{\mathbb{C}}) .
\]
\struck{\ill{}}

Considérons le noyau \ill{} qui contient le noyau \ill{} de
$\mathrm{Gal}(k/k_{0}) \to \mathrm{Gal}(k/\bar{k}_{0})$. Ce noyau est envoyé
\ill{} dans $G(\mathbb{A}_{\mathbb{C}})$ \ill{} par $\varphi_{v,\infty}$, et
l'image \ill{} \struck{Il en résulte que \struck{$\varphi$} \ill{} plus précis},
\ill{}

\ill{} \emph{dominant} $\varphi_{v}\bigl(\mathrm{Gal}(k/k_{0})\bigr)$ \ill{} un
\ill{} (ouvert) \ill{} de $G(\mathbb{A}_{\mathbb{C}})$.

\note{La phrase est coupée en bas du feuillet : la page 61, premier feuillet du
lot suivant, la reprend et écrit la relation $M^{\mathbb{Z}}_{w} = g
M^{\mathbb{Z}}_{v}$ pour $w = \sigma v$ ainsi que le cocycle conjugué
$\varphi_{w}(\sigma) = g_{\infty} \varphi_{v}(\sigma) g_{\infty}^{-1}$. Le bas
du feuillet porte encore, à la mine, « ( )~? ».}

\end{document}
